% Bagian: sets-functions-relations
% Bab: size-of-sets
% Subbagian: comparing-sizes

\documentclass[../../../include/open-logic-section]{subfiles}

\begin{document}

\olfileid[id]{sfr}{siz}{car}

\olsection{Himpunan dengan Ukuran Berbeda dan Teorema Cantor}

\begin{explain}
Kita telah memberikan pernyataan saksama bagi gagasan bahwa dua himpunan
berukuran sama. Kita juga dapat memberikan pernyataan saksama bagi gagasan
bahwa satu himpunan lebih kecil daripada himpunan lain. Definisi kita tentang
``lebih kecil daripada (atau berkardinalitas sama)'' tidak memerlukan
!!a{bijection} antara kedua himpunan, melainkan !!a{injection} dari himpunan
pertama ke himpunan kedua. Jika fungsi semacam itu ada, ukuran himpunan pertama
kurang dari atau sama dengan ukuran himpunan kedua. Secara intuitif,
!!a{injection} dari satu himpunan ke himpunan lain menjamin bahwa daerah hasil
fungsi mempunyai sekurang-kurangnya sebanyak !!{element}s domainnya, karena
tidak ada dua !!{element}s domain yang dipetakan ke !!{element} daerah hasil
yang sama.
\end{explain}

\begin{defn}
$A$ \emph{tidak lebih besar daripada}~$B$, ditulis $\cardle{A}{B}$, jika dan
hanya jika terdapat !!a{injection} $f \colon A \to B$.
\end{defn}

Jelas bahwa relasi ini refleksif dan transitif, tetapi tidak simetris (hal ini
ditinggalkan sebagai latihan). Kita juga dapat memperkenalkan suatu pengertian
yang menyatakan bahwa satu himpunan (secara ketat) lebih kecil daripada
himpunan lain.

\begin{defn}
$A$ \emph{lebih kecil daripada}~$B$, ditulis $\cardless{A}{B}$, jika dan hanya
jika terdapat !!a{injection}~$f\colon A \to B$, tetapi tidak terdapat
!!{bijection}~$g\colon A \to B$; yakni $\cardle{A}{B}$ dan
$\cardneq{A}{B}$.
\end{defn}

Jelas bahwa relasi ini irefleksif dan transitif. (Hal ini ditinggalkan sebagai
latihan.) Dengan menggunakan notasi ini, kita dapat mengatakan bahwa suatu
himpunan $A$ adalah !!{enumerable} jika dan hanya jika $\cardle{A}{\Nat}$,
serta bahwa $A$ adalah !!{nonenumerable} jika dan hanya jika
$\cardless{\Nat}{A}$. Ini memungkinkan kita menyatakan kembali
\oliflabeldef{sfr:siz:nen-alt:thm:nonenum-pownat}{%
\olref[sfr][siz][nen-alt]{thm:nonenum-pownat}
sebagai pengamatan bahwa
$\cardless{\Nat}{\Pow{\Nat}}$}{%
\olref[sfr][siz][nen]{thm:nonenum-pownat}
sebagai pengamatan bahwa $\cardless{\PosInt}{\Pow{\PosInt}}$}. Bahkan,
\citet{Cantor1892} membuktikan bahwa butir terakhir ini berlaku \emph{secara
sepenuhnya umum}:

\begin{thm}[Cantor]\ollabel{thm:cantor}
$\cardless{A}{\Pow{A}}$, untuk setiap himpunan $A$.
\end{thm}

\begin{proof}
Pemetaan $f(x) = \{x\}$ adalah !!a{injection} $f \colon A \to \Pow{A}$,
karena jika $x \neq y$, maka $\{x\} \neq \{y\}$ juga menurut ekstensionalitas,
sehingga $f(x) \neq f(y)$. Jadi, kita memperoleh $\cardle{A}{\Pow{A}}$.

\begin{editorial}
Kita menyajikan bukti yang lebih perlahan jika \olref[nen]{sec} tersedia;
jika tidak, kita menyajikan bukti lebih cepat yang selaras dengan
\olref[nen-alt]{sec}.
\end{editorial}

\oliflabeldef{sfr:siz:nen:sec}{%
  Sekarang kita akan menunjukkan bahwa tidak mungkin ada fungsi
  !!a{surjective}~$g\colon A \to \Pow{A}$, apalagi fungsi !!a{bijective}, dan
  dengan demikian $\cardneq{A}{\Pow{A}}$. Andaikan $g\colon A \to \Pow{A}$.
  Karena $g$ total, setiap $x \in A$ dipetakan ke himpunan bagian $g(x)
  \subseteq A$. Kita dapat menunjukkan bahwa $g$ tidak mungkin surjektif.
  Untuk melakukannya, kita mendefinisikan himpunan bagian~$\overline{A}
  \subseteq A$ yang, menurut definisinya, tidak mungkin berada dalam daerah
  hasil~$g$. Misalkan
  \[
  \overline{A} = \Setabs{x \in A}{x \notin g(x)}.
  \]
  Karena $g(x)$ terdefinisi untuk semua $x \in A$, $\overline{A}$ jelas
  merupakan himpunan bagian~$A$ yang terdefinisi dengan baik. Namun,
  himpunan itu tidak mungkin berada dalam daerah hasil~$g$. Ambil sebarang
  $x \in A$; kita akan menunjukkan bahwa $\overline{A} \neq g(x)$. Jika $x
  \in g(x)$, maka ia tidak memenuhi $x \notin g(x)$, sehingga menurut
  definisi~$\overline{A}$, berlaku $x \notin \overline{A}$. Jika $x \in
  \overline{A}$, ia harus memenuhi sifat pendefinisi
  ~$\overline{A}$, yakni $x \in A$ dan $x \notin g(x)$. Karena $x$ dipilih
  sebarang, hal ini menunjukkan bahwa untuk setiap $x \in A$, $x \in g(x)$
  jika dan hanya jika $x \notin \overline{A}$, sehingga $g(x) \neq
  \overline{A}$. Dengan kata lain, $\overline{A}$ tidak mungkin berada dalam
  daerah hasil~$g$, bertentangan dengan pengandaian bahwa~$g$ surjektif.}{
  Tinggal ditunjukkan bahwa $\cardneq{A}{\Pow{A}}$. Untuk memperoleh
  kontradiksi, andaikan $\cardeq{A}{\Pow{A}}$, yakni terdapat suatu
  !!{bijection} $g \colon A \to \Pow{A}$. Sekarang perhatikan:
  \[
    D = \Setabs{x \in A}{x \notin g(x)}
  \]
  Perhatikan bahwa $D \subseteq A$, sehingga $D \in \Pow{A}$. Karena $g$
  adalah !!a{bijection}, terdapat suatu $y \in A$ sedemikian sehingga $g(y) =
  D$. Namun, sekarang kita memperoleh:
  \[
    y \in g(y) \text{ jika dan hanya jika } y \in D \text{ jika dan hanya jika } y \notin g(y).
  \]
  Ini merupakan kontradiksi; jadi $\cardneq{A}{\Pow{A}}$.}{}
\end{proof}

\begin{explain}
\oliflabeldef{sfr:siz:nen:thm:nonenum-pownat}{Akan bermanfaat jika kita
  membandingkan bukti \olref{thm:cantor} dengan bukti
  \olref[nen]{thm:nonenum-pownat}. Di sana kita menunjukkan bahwa dari setiap
  daftar $Z_1$, $Z_2$, \dots, dari himpunan bagian~$\PosInt$, dapat dibangun
  suatu himpunan bilangan~$\overline{Z}$ yang dijamin tidak berada dalam
  daftar. Jaminan itu diperoleh karena, untuk setiap $n \in \PosInt$, berlaku
  $n \in Z_n$ jika dan hanya jika $n \notin \overline{Z}$. Dengan demikian,
  selalu terdapat suatu bilangan yang merupakan !!a{element} dari salah satu
  di antara $Z_n$ dan $\overline{Z}$, tetapi bukan yang lain. Kita mengikuti
  gagasan yang sama di sini, hanya saja indeks~$n$ kini merupakan !!{element}s
  dari~$A$, bukan dari~$\PosInt$. Himpunan $\overline{A}$ didefinisikan agar
  berbeda dari~$g(x)$ untuk setiap $x \in A$, karena $x \in g(x)$ jika dan
  hanya jika $x \notin \overline{A}$. Sekali lagi, selalu ada !!a{element}
  dari~$A$ yang merupakan !!a{element} dari salah satu di antara $g(x)$ dan
  $\overline{A}$, tetapi bukan yang lain. Sebagaimana $\overline{Z}$ tidak
  mungkin berada dalam daftar $Z_1$, $Z_2$, \dots, $\overline{A}$ juga tidak
  mungkin berada dalam daerah hasil~$g$.}{}

\oliflabeldef{sfr:siz:nen-alt:thm:nonenum-pownat}{Akan bermanfaat jika kita
  membandingkan bukti \olref{thm:cantor} dengan bukti
  \olref[nen-alt]{thm:nonenum-pownat}. Di sana kita menunjukkan bahwa dari
  setiap daftar $N_0$, $N_1$, $N_2$, \dots, dari himpunan bagian~$\Nat$, kita
  dapat membangun suatu himpunan bilangan~$D$ yang dijamin tidak berada dalam
  daftar. Jaminan itu diperoleh karena $n \in N_n$ jika dan hanya jika $n
  \notin D$, untuk setiap $n \in \Nat$. Kita mengikuti gagasan yang sama di
  sini, hanya saja indeks~$n$ kini merupakan !!{element}s dari~$A$, bukan dari
  ~$\Nat$. Himpunan $D$ didefinisikan agar berbeda dari~$g(x)$ untuk setiap
  $x \in A$, karena $x \in g(x)$ jika dan hanya jika $x \notin D$.}{}

Bukti tersebut juga patut dibandingkan dengan bukti Paradoks Russell,
\olref[sfr][set][rus]{thm:russells-paradox}. Memang, Teorema Cantor menjadi
inspirasi bagi paradoks Russell sendiri.
\end{explain}

\begin{prob}
  Tunjukkan bahwa tidak mungkin ada !!a{injection} $g\colon \Pow{A} \to A$,
  untuk setiap himpunan~$A$. Petunjuk: Andaikan $g\colon \Pow{A} \to A$
  !!{injective}. Perhatikan $D = \Setabs{g(B)}{B \subseteq A \text{ dan }
  g(B) \notin B}$. Misalkan $x = g(D)$. Gunakan fakta bahwa $g$ !!{injective}
  untuk memperoleh kontradiksi.
\end{prob}

\end{document}
